\documentclass[conference]{IEEEtran}
\IEEEoverridecommandlockouts
\usepackage{cite}
\usepackage{amsmath,amssymb,amsfonts}
\usepackage{algorithmic}
\usepackage{graphicx}
\usepackage{textcomp}
\usepackage{xcolor}
\usepackage{booktabs}
\def\BibTeX{{\rm B\kern-.05em{\sc i\kern-.025em b}\kern-.08em
    T\kern-.1667em\lower.7ex\hbox{E}\kern-.125emX}}
\begin{document}

\title{Prompt Structure and Human Preference in Generative Image Models}

\author{\IEEEauthorblockN{Muziyan Wang}
\IEEEauthorblockA{\textit{Dept. of Economics} \\
\textit{University of Michigan}\\
Ann Arbor, MI, USA \\
muziyanw@umich.edu}
}

\maketitle

\begin{abstract}
This project investigates how prompt-level structural and emotional features affect
human preference scores for images generated by large text-to-image model (GPT-4o).
Using a subset of dataset from Hugging Face and a range of prompt features,
we train a classifier to predict user preferences and explore
which prompt elements most influence preference formation. Our findings
shed light on effective prompt engineering strategies and contribute to
human-centric evaluation methods in generative AI.
\end{abstract}

\begin{IEEEkeywords}
prompt structure, image generation, human preference, emotional lexicon, logistic regression
\end{IEEEkeywords}

\section{Introduction}
As generative AI systems become more prevalent in visual content creation, understanding how users perceive the quality of generated images becomes crucial. Among the most influential factors is the language used in text prompts. Prior work has focused on architectural improvements or perceptual metrics, but the role of the prompt itself in shaping preference remains underexplored. This study aims to fill this gap by identifying structural and emotional prompt features that correlate with preference for GPT-4o-generated images.

We use the NRC Emotion Lexicon~\cite{nrc} to extract affective features such as emotion word presence and category diversity.
Recent work such as TIFA~\cite{tifa} and PromptBench~\cite{promptbench} has emphasized the importance of prompt content in shaping human preferences for AI-generated outputs. TIFA benchmarks vision-language alignment using textual feedback, while PromptBench evaluates prompts based on fluency, informativeness, and relevance.
To better capture subjective judgments, ImageReward~\cite{imagereward} introduces a reward model trained on 137k expert comparisons, improving generation quality by modeling explicit user preferences.
Finally, SHAP~\cite{shap} provides interpretable explanations for model decisions, highlighting the importance of emotional and structural prompt characteristics.

\section{Dataset and Feature Construction}
Due to memory limitations on the local machine, only a subset of the full OpenAI-4o Human Preference dataset available on Hugging Face was downloaded. The subset contains image-prompt pairs and user preference comparisons between two generations. To operationalize the classification task, we labeled image 2 as preferred (1) or not preferred (0) based on the provided comparison field.

We extract prompt-level features for each record, including:
\begin{itemize}
\item Prompt length (word count)
\item Number of emotional words (using NRC Emotion Lexicon)
\item Number of distinct emotional categories
\item Presence of role-setting expressions
\item Structural constraints (e.g., "must include")
\item Presence of color-related words
\item Coherence and alignment scores associated with image 2
\end{itemize}

These features were chosen based on linguistic theory, prior literature, and preliminary exploration of the dataset.

\section{Methodology}
Our approach follows a four-step pipeline: feature engineering, exploratory data analysis, predictive modeling, and insight generation.

\textbf{Step 1:} Feature Annotation and Exploratory Data Analysis (EDA)
We annotated the prompt data with structural and emotional features, including prompt length, emotional word count, emotional category diversity, narrative framing, and color presence. These features were engineered using regular expressions and the NRC Emotion Lexicon.

Figure~\ref{fig:heatmap} presents a correlation heatmap among all prompt-level features. Emotional word count demonstrates the strongest positive correlation with preference label (\textit{r} = 0.47), reinforcing its predictive value. In contrast, model-derived alignment and coherence scores exhibit weak correlations with preference, motivating our emphasis on linguistic cues.

\begin{figure}[htbp]
\centerline{\includegraphics[width=0.48\textwidth]{eda_heatmap_correlations.png}}
\caption{Correlation heatmap of prompt-level features}
\label{fig:heatmap}
\end{figure}

\textbf{Step 2–4: Modeling and Evaluation}
The remaining steps involved predictive modeling using logistic regression, followed by SHAP-based feature attribution and visual interpretation.

To evaluate model performance, we use standard classification metrics including accuracy, precision, recall, and F1-score. We analyze results from logistic regression using a consistent 80/20 train-test split with a fixed random seed. Additionally, SHAP values are computed to assess feature-level impact on predictions.

We used logistic regression to model the probability that a given prompt receives higher human preference. The model estimates the log-odds as a linear combination of engineered features:

\begin{equation}
P(y=1 \mid \mathbf{x}) = \frac{1}{1 + e^{-(\beta_0 + \beta_1 x_1 + \beta_2 x_2 + \cdots + \beta_p x_p)}}
\label{eq:logistic}
\end{equation}

The model was trained using binary cross-entropy loss on the preference labels, and coefficients were later analyzed using SHAP for interpretability.

\section{Results}
The model achieved an overall accuracy of 80\% on the test set. The performance metrics are summarized in Table~\ref{tab:classification}.

\begin{table}[htbp]
\caption{Classification Report Summary}
\label{tab:classification}
\centering
\begin{tabular}{@{}lcccc@{}}
\toprule
Class & Precision & Recall & F1-score & Support \\
\midrule
0 (not preferred) & 0.81 & 0.83 & 0.82 & 115 \\
1 (preferred)     & 0.77 & 0.74 & 0.75 & 85 \\
\midrule
\textbf{Accuracy} & \multicolumn{4}{c}{\textbf{0.80}} \\
\bottomrule
\end{tabular}
\end{table}

\noindent Confusion matrix:

$$
\begin{bmatrix}
96 & 19 \\
22 & 63
\end{bmatrix}
$$

\section{Feature Importance}
To better understand what features influenced model predictions, we computed SHAP values using the trained model. Figure~\ref{fig:shap} shows that emotional word count, emotional category variety, and presence of color words were among the most influential predictors of preference.

\begin{figure}[htbp]
\centerline{\includegraphics[width=0.48\textwidth]{shap_nrc_color.png}}
\caption{SHAP value summary plot}
\label{fig:shap}
\end{figure}

These findings validate our hypothesis that emotional and structural signals in prompts play a key role in guiding human judgment, even before image content is considered.

\section{Conclusion}
This study demonstrates that human preference in text-to-image generation can be effectively predicted using prompt-level features alone. Emotionally expressive language, role-setting narrative forms, and the presence of color terms emerged as strong predictors of preferred outputs.

While the current study is based on a partial sample of the OpenAI-4o Human Preference dataset, the results show that even with limited access to data, consistent patterns of human judgment can be modeled using lightweight, interpretable features. We also highlight the utility of explainable models such as logistic regression and SHAP in guiding prompt engineering strategies and advancing user-centric AI evaluation frameworks.

Future work may include extending this framework with semantic embeddings (e.g., BERT or CLIP), validating results with larger-scale annotations, and conducting A/B testing to derive practical prompt optimization guidelines.

\section*{Acknowledgment}
This research project was developed for the final project of STATS 507 at the University of Michigan. The author thanks the instructional team for their guidance.

\bibliographystyle{IEEEtran}
\begin{thebibliography}{1}

    \bibitem{nrc}
    S. Mohammad and P. Turney, "NRC Emotion Lexicon," National Research Council Canada, 2013.
    \bibitem{tifa}
    P. Wu et al., “TIFA: An Evaluation Benchmark for Multi-modal Language Understanding via Text-to-Image Feedback,” arXiv:2305.13082, 2023.
    \bibitem{promptbench}
    Z. Liu et al., “PromptBench: Towards Evaluating the Quality of AI-Generated Prompts,” arXiv:2311.16417, 2023.
    \bibitem{imagereward}
    J. Xu, X. Liu, Y. Wu, Y. Tong, Q. Li, M. Ding, J. Tang, and Y. Dong, “ImageReward: Learning and Evaluating Human Preferences for Text-to-Image Generation,” in Advances in Neural Information Processing Systems (NeurIPS), 2023.
    \bibitem{shap}
    S. M. Lundberg and S.-I. Lee, “A unified approach to interpreting model predictions,” in Advances in Neural Information Processing Systems (NeurIPS), vol. 30, pp. 4765–4774, 2017.
    \end{thebibliography}


\end{document}

